\section*{Appendix A: External tools actually used (n=29)}
\addcontentsline{toc}{section}{Appendix A}
Each row is an executed use; evidence files are cited in
\texttt{results/EXTERNAL\_TOOLS.md} and the individual
\texttt{results/*.json} files. Tools listed as planned/staged in the
ledger but not executed are NOT counted.
\begin{longtable}{r p{4.5cm} p{9.5cm}}
\hline
\# & Tool & What it was used for \\ \hline
\endhead
1 & RCSB PDB / data.rcsb.org & 481 WT+mutant structures downloaded, parsed (graph construction) \\
2 & ProtDDG-Bench (github.com/protddg-bench) & S2648/Ssym/P53 train-test protocol \\
3 & Pucci et al. 2018 Bioinformatics (bty348) published Table 1 & benchmark reference values (PoPMuSiCsym 0.48/1.62 etc.), verified from PDF \\
4 & FireProtDB 2.0 REST API & SOD1 mutation search + mutant detail endpoints \\
5 & ThermoMutDB API & full dump, human SOD1 (P00441) ddG extraction \\
6 & Europe PMC REST API & article lookup (PMID/PMC resolution) \\
7 & Zenodo API & ThermoMPNN FireProtDB+PDB dataset (record 8169289) \\
8 & NumPy & all model math (GCN layers, Adam, backprop) \\
9 & pytest & 24 hermetic unit tests \\
10 & git + GitHub & versioned delivery (SSH, ed25519) \\
11 & Google Drive & results delivery folder (bundle uploads) \\
12 & SoDCoD (fujisawagroup) & SOD1 mutant conformation metadata (lookup) \\
13 & FireProtDB OpenAPI/Swagger spec & API query construction \\
14 & matplotlib & paper figures (fig1 scatter, fig2 bias bars, fig3 curves) \\
15 & pdflatex / TeX Live (mathptmx) & paper PDF compilation \\
16 & Biopython (Bio.PDB) & independent cross-check of our PDB parser; exposed the altloc-CA duplicate bug (2VUK 197 vs 195 residues); confirms 195-residue chain + superstable background in stats\_and\_baseline \\
17 & SciPy (scipy.stats) & Spearman rho/p-value, Fisher-z CI on benchmark r \\
18 & scikit-learn (Ridge) & identity-only no-structure baseline on same train/test rows (r\_inv 0.251 vs GNN 0.592) \\
19 & UniProt REST API & P04637/P00441 feature cross-check: all 11 p53 scan positions carry curated cancer-variant annotations; R273/Y220 are LFS germline hotspots; SOD1 H46R/G85R/V148I ALS1-annotated (res \\
20 & PDBe API (SIFTS) & PDB<->UniProt numbering verification: 2VUK and 1SPD author numbering = UniProt canonical at all 10 checked positions (results/pdbe\_sifts\_check.json); guards the SOD1 numbering-ambi \\
21 & FreeSASA 2.2.1 & per-residue SASA on 2VUK; burial correlates with predicted |ddG| (Spearman 0.391, p=1.6e-8) (results/burial\_analysis\_2vuk.json) \\
22 & pandas & scan/candidate table joins and exports (results/scan\_candidates\_2vuk\_top250.csv, burial analysis) \\
23 & AlphaFold Protein Structure DB & AF-P00441-F1 model downloaded; pLDDT-vs-scan analysis: top SOD1 candidates sit in high-confidence regions (pLDDT ~98.8); pLDDT/|ddG| Spearman 0.295 (results/alphafold\_db\_lookup.jso \\
24 & Pfam via EBI InterPro API & domain membership: all 11 p53 scan positions inside PF00870 (P53 DNA-binding); SOD1 positions inside PF00080 except A4/C6 (N-terminal, outside domain - coincides with the two worst \\
25 & RCSB PDB Data API (entry annotations) & experimental metadata: 2VUK 1.5A Y220C+stabilizing-drug complex, 1SPD 2.4A, 3ECU apo 1.9A, 1N18 C6A/C111S 2.0A (results/rcsb\_entry\_metadata.json) \\
26 & NCBI E-utilities (PubMed esearch/esummary) & literature verification for A4V apo-stability caveat (3 hits, titles committed); p53 suppressor phrase queries returned 0 hits (recorded honestly) (results/ncbi\_eutils\_litcheck.jso \\
27 & gnomAD GraphQL API & TP53 germline constraint: pLI 0.9996, mis\_z 1.12, oe\_lof 0.258 (results/gnomad\_tp53\_constraint.json) \\
28 & PyMOL open-source 3.x & structure figure fig4 (2VUK candidate geometry, paper/figs/fig4\_candidates\_2vuk.png; session build/pymol\_2vuk\_session.pse) \\
29 & STRING DB API v12 & TP53/SOD1 functional-network context (results/string\_network\_context.json) \\
\hline
\end{longtable}
